\section{Introduction}
\label{ch:intro}

\subsection{Starting from Ill-Posedness}

Many inference tasks are ill-posed at their root. The brightness of a single pixel can correspond to infinitely many depths, and a single measurement to infinitely many states; no local observation alone can ever fix the answer. Take the most ordinary question in real-time 3D reconstruction: given one frame, how far is the surface at some pixel from the camera? The color of the pixel can be weak light reflected from a nearby dark surface, or strong light reflected from a distant bright surface, and under transparency, occlusion, or repetitive texture it can correspond to still more complex structures. The relation between observation and structure is not one-to-one but many-to-many, and this is exactly ill-posedness in Hadamard's sense: the solution does not exist, is not unique, or does not depend continuously on the observation \citep{hadamard1902}.

Ill-posedness is not caused by engineering crudeness; it comes from the task itself. An observation is the residue of structure after projection, sampling, and noise contamination; information has already been lost in the forward process, and no local operation in the inverse direction can bring the lost information back. Early vision theory stated this systematically: stereo vision, motion estimation, optical flow, and shape reconstruction are all ill-posed problems that admit a unique solution only after additional constraints are introduced \citep{poggio1985}. These additional constraints are what this paper calls context: neighboring pixels usually belong to the same surface, neighboring frames usually correspond to the same motion, and neighboring voxels usually lie within the same organ. Context pins down the degrees of freedom that local observations cannot fix.

The trouble is that context cannot be applied uniformly. Inside a surface the constraint should be strong, so that the neighborhood corroborates itself and noise is smoothed away; at boundaries and occlusions it should be cut, so that the two sides take shape independently, otherwise depth blurs into a smear along object edges. When motion is smooth, consecutive frames can corroborate each other; when motion is abrupt, they must be independent, otherwise the estimate is dragged along by history. In other words, how strongly context should be used is itself a quantity that varies with position and time, and its correct value depends on where the boundaries are in the current scene, where the motion breaks, and along which contour the occlusion occurs. This property of variable strength that changes with position and time is exactly the most intractable part of this class of tasks: a context rule is not a fixed and unchanging law but a stipulation that varies with the evidence.

\subsection{Open Inversion Tasks}

Organizing the observations above into four items yields the object of study of this paper. First, observations are ill-posed: local observations do not suffice to uniquely determine the structure being sought, so context is indispensable. Second, the required context strength is heterogeneous: the strength needed varies across positions in space and across instants in time, and where it should be strong and where weak depends on the content of the current evidence; that is, the locations of boundaries are content-addressed and cannot be tabulated in advance. Third, inputs at inference time may come from outside the training distribution: the scenes, sensors, and motion patterns of the deployment environment are not bounded by the coverage of the training corpus. This paper calls the set of such inputs \emph{the open set}, and the property that inference inputs lie in it \emph{distributional openness}. Fourth, the demand is non-freezable: the required context rule, as a mapping from evidence to strengths, cannot be given within any pre-given tolerance on the open set by any values frozen before inference begins. This paper calls an inference task satisfying all four items an \emph{open inversion task}.

Real-time 3D reconstruction is its most typical representative: the depth field is ill-posed pixel by pixel, the smoothing strength must be heterogeneous inside and outside boundaries, and the environments that a robot or camera enters are not constrained by the training set; recent real-time radiance-field reconstruction has pushed this form to the engineering foreground \citep{kerbl2023}. Image and video restoration belong to the same class: the mixture of degradation kernels and noise levels makes local observations ill-posed, textured regions and flat regions demand different restoration strengths, and real degradations exceed the coverage of synthetic training sets. Online state estimation is no different: the measurement equation is ill-posed along weakly observable directions, the process-noise strength switches with the motion mode, and the open world keeps supplying out-of-distribution operating conditions. Medical and scientific inversion, from tomographic reconstruction to inverse scattering, also fall within this framework.

This paper takes the four items as premises rather than conclusions. They hold on large classes of real tasks: inverse problems are mostly ill-posed, which is the defining feature of Hadamard's criterion; natural data are mostly heterogeneous, since boundaries, occlusions, and abrupt changes are the structure of natural signals rather than exceptions; online systems are mostly open, and statistical learning theory long ago pointed out that once the assumption of agreement between the training and test distributions is removed, the guarantees based on distribution agreement lapse with it \citep{vapnik1995}. But the four items are not universally true. Globally homogeneous interpolation tasks fail heterogeneity, because the required smoothing strength is the same everywhere; in fixed-grid numerical simulation, boundary positions are set by the task rather than by the content of evidence; in state estimation under fixed operating conditions, the strength demand is given by an exogenous schedule. These three kinds of tasks are not open inversion tasks, and fixed-memory components on them are not among what this paper falsifies; this exactly delimits the applicable boundary of the conclusion. The boundary of the fourth item is equally concrete: if the demand can be exactly identified from a fixed corpus, for example when the demand is a simple linear function of the evidence that a few noiseless samples pin down, a frozen component can be complete on such tasks, which are therefore not open inversion tasks; the falsification stops exactly where the demand becomes freezable. In the terms of Section~\ref{sec:p4}, such a task fails the barrier condition: its demand has a pre-given modulus, and no separation survives at fine scale.

\subsection{The Claim to Be Falsified}

Since context rules are the core of this class of tasks, a natural question arises: does there exist a component that, by itself alone, applies context rules correctly on every open inversion task? This question is no straw man; it has a definite real-world form in public discourse, as the conjunction of two claims. The first is the engineering corollary of universal approximation theorems \citep{cybenko1989,hornik1989}: since neural networks can approximate arbitrary functions, and inference is likewise a mapping from observations to structures, a sufficiently large network with a sufficiently large corpus should in principle suffice to carry out inference; the systems literature on end-to-end 3D reconstruction and online estimation is generally built on this implicit premise. The second is the scale claim: deficiencies in contextual ability are regarded as a matter of corpus and parameter scale rather than of architectural category, and the capability discourse of the large-model era has pushed this claim to its peak \citep{bubeck2023}, while its critics point out that some so-called emergent abilities may be artifacts of metric choice \citep{schaeffer2023}. Taken together, the two claims assert that there exists some fixed-memory component which, once at scale, completes the task on its own. What this paper falsifies is exactly this claim:

\begin{claim}[$\mathcal{R}$]
There exists a fixed-memory component that, using only that component and without any external proposal or arbitration stage, is complete for open inversion tasks; that is, it completes the task independently in the sense of the open set, of no distribution agreement, and of cross-scenario reuse.
\end{claim}

A fixed-memory component is one whose context rules are frozen into parameters that do not change at inference time and whose parameter values come from training-corpus statistics or manual setting. Its counterpart in the existing literature is the carrier of parametric amortized inference: inference is performed by a forward process with fixed parameters whose values are determined by the corpus. It should be noted that the Mode A side of this paper does not exclude parameters; analytic invariants are parameters too, and the difference lies in the source being calibration and geometry rather than corpus statistics. Completeness in this paper means independent completion in the sense of task capability and has nothing to do with the completeness of formal systems in mathematical logic. Neural networks are the canonical instance of fixed-memory components: they carry the context rule in weights obtained from training or calibration. The instantiation of Claim $\mathcal{R}$ on this carrier is the neural-network-only claim $\mathcal{P}$, and Theorem 5 of this paper is exactly this instantiation.

The qualifier ``only'' is crucial. What is falsified is solo use, not the usefulness of these components within combinations. Neural networks as feature extractors can play roles in larger systems, and this paper denies none of that. One misreading must also be ruled out: achieving high accuracy on a fixed benchmark close to the training distribution is within reach of both kinds of networks, but that is a \emph{weak claim} and is not the object of this paper. The question asked here is whether the component alone, as a general-purpose inference mechanism, can complete the task independently, that is, whether it possesses completeness in the task sense. The truth of the weak claim and the falsity of the strong claim are compatible, just as an instrument of extreme accuracy under standard operating conditions may still fail to work independently under arbitrary operating conditions.

\subsection{Method of Falsification and Structure of the Paper}

The methodology of this paper is a deductive impossibility argument under explicit criteria. Claim $\mathcal{R}$ is an existential claim, and its negation is a universal statement about an entire task class, so the Popperian shortcut of refuting a universal claim by a single counterexample \citep{popper1959} is not available here: negating an existential claim does not require running through all tasks, but it does require a deductive chain from the defining criteria of the task class to the capability boundary of the component, and the criteria must pin down the task class first. All derivations in this paper are formal and rely on no experimental data: experiments can exhibit symptoms, but what this paper proves is the crux. One point about the shape of the argument deserves mention here, because it decides how far the negative result reaches. The algebraic core does not depend on the quadratic smoothing form of the classical variational problem: Lemma 1 is proved for a separable data term with an arbitrary even differentiable pairwise potential, and no convexity is used, so the exclusion covers the stationary points of nonconvex objectives as well. What it excludes is the frozen coefficient family in general, with the global scalar as its shallowest case, and the minimal obstruction is stated as an equivalence rather than as a sufficient condition: one single position whose required strength ratio changes between two admissible configurations excludes every frozen family, whatever number of entries it has. Whether the criteria capture real tasks is an empirical question, and this paper discusses the conditions under which each criterion holds; what lies outside the criteria is not a counterexample but a boundary of the class. The word \emph{falsification} is used in its deductive sense throughout: what is refuted is Claim $\mathcal{R}$ under Criteria P1 to P4, by deduction rather than by experiment, and Section~\ref{ch:verify} delimits the term once more where the machine verification is described.

